\section{Legal Landscape}
\label{sec:legal}

\subsection{Explicit Statement of Legal Non-Adoption}

\textbf{Declination has been accepted as evidence by no federal or state court as
of 2026.} This is the central legal fact that expert witnesses must acknowledge
when discussing this metric. Unlike EG (extensively discussed in \textit{Gill}
and \textit{Rucho}), MM (cited in \textit{Harper v.\ Hall}), and Partisan Bias
(cited in \textit{Harper v.\ Hall} and PA proceedings), Declination has not
been credited by any court as a basis for any redistricting decision.

The legal record is as follows:

\textbf{Common Cause v.\ Rucho (M.D.N.C.\ 2018).}\footnote{\textit{Common Cause
v.\ Rucho}, 318 F.Supp.3d 777 (M.D.N.C.\ 2018).} Declination was introduced as
an exhibit in the North Carolina district court proceedings that produced one of
the two cases consolidated before the Supreme Court in \textit{Rucho}. However,
the district court's opinion---which ruled for the plaintiffs on partisan
gerrymandering grounds---did not rely on Declination. The court's analysis
credited EG, MM, and Partisan Bias evidence; Declination was presented in expert
testimony but was not cited in the court's findings. When SCOTUS reversed on
political question grounds in \textit{Rucho}, the Declination evidence was
not part of the appellate record that received any attention.

\textbf{Harper v.\ Hall (N.C.\ 2022).} The North Carolina Supreme Court's opinion
invalidating the 2022 congressional map cited EG and MM as metrics demonstrating
systematic partisan advantage. The opinion did not mention Declination. North
Carolina's redistricting litigation proceedings are among the most thoroughly
documented state court redistricting cases in the country; the absence of
Declination from the court's analysis is a considered omission, not an oversight.

\textbf{Wisconsin, Pennsylvania, Ohio proceedings.} None of the active state court
redistricting proceedings through 2026 in these three jurisdictions has cited
Declination in a court opinion.

\subsection{Why Declination Has Not Been Adopted}

The three methodological critiques in Section~\ref{sec:definition} provide the
likely explanation for judicial non-adoption:

\textbf{The undefined-when-swept problem} is disqualifying for small delegations
in wave elections and for historically lopsided states. A metric that cannot be
computed in some scenarios provides an unreliable foundation for a constitutional
standard. Courts developing redistricting doctrine across the full range of states
and election environments cannot rely on a metric that fails for Massachusetts in
a strong Democratic year.

\textbf{Threshold sensitivity} produces discontinuous metric values when districts
are near 50\%---exactly the scenario most relevant in contested redistricting cases.
A metric that jumps discontinuously when a district's vote share crosses 50\% is
vulnerable to the criticism that it produces different answers depending on turnout
variation in individual precincts, not the structural properties of the map.

\textbf{The absence of a calibrated threshold} means courts have no basis for
determining when $\delta > 0.20$ or $\delta > 0.30$ indicates a constitutional
violation versus permissible partisanship. Without calibration to historical
data on partisan outcomes under neutral redistricting, Declination cannot be
operationalized as a constitutional standard.

\subsection{Future Prospects and Expert Witness Guidance}

Declination's geometric transparency---the ability to read it directly from
a plot of sorted district vote shares---gives it potential as a courtroom
exhibit if its methodological critiques can be addressed. A hybrid exhibit
that shows the sorted vote-share plot (the visual intuition of Declination)
alongside EG and MM (the legally credentialed metrics) might communicate the
packing-cracking geometry more effectively than any of the scalar metrics alone.

Expert witnesses using Declination should:
\begin{enumerate}
\item Lead with the visual exhibit (sorted vote-share plot) rather than the
  scalar $\delta$ value, letting the geometric asymmetry speak before introducing
  the formal metric.
\item Explicitly disclose that no court has adopted Declination as a standard
  and that the metric's legal status is academic as of 2026.
\item Acknowledge and respond to the three methodological critiques.
\item Present $\delta$ alongside EG, MM, and Partisan Bias to demonstrate
  that the multi-metric conclusion is consistent regardless of whether
  Declination is credited as a legal standard.
\end{enumerate}
